% --- Hasil dan Pembahasan ---
\chapter{Hasil Penelitian dan Pembahasan}

%=========================================================%
%          TULIS ISI HASIL DAN PEMBAHASAN DI SINI         %
% Jangan Lupa untuk Menghapus Contoh Tulisan di Bawah Ini %
%          (Termasuk Panduan di Setiap Subbabnya)         %
% Contoh Teks Dapat Di-comment dengan Tanda "%" di Depan  %
%   jika Ingin Mempertahankan Panduan di Dalam File Ini   %
%=========================================================%

%------------------------ IMBAUAN ------------------------%
% Di bawah ini menggunakan tiga subbab, yang disarankan   %
% untuk memisah gambaran umum, penyajian data, dan        %
% bahasan. Jika pembimbing Anda kurang menyukai ini, Anda %
% dapat membuang subbab "Penyajian Hasil Penelitian" di   %
% bawah ini.                                              %
%---------------------------------------------------------%

% --- Gambaran Umum ---
\section{Gambaran Umum Subjek dan Objek Penelitian}

Jelaskan secara ringkas karakteristik dari objek penelitian Anda (misalnya profil perusahaan, latar belakang instansi, atau deskripsi sistem yang diamati). Jika penelitian Anda menggunakan kuesioner atau wawancara, sajikan profil demografi responden secara agregat (seperti distribusi usia, jenis kelamin, tingkat jabatan, atau latar belakang pendidikan) menggunakan tabel atau grafik yang rapi. Laporkan juga bagaimana proses pelaksanaan penelitian di lapangan secara kronologis.



% --- Penyajian Hasil ---
\section{Penyajian Hasil Penelitian}

Sajikan data-data temuan ilmiah yang telah Anda kumpulkan dan olah dari lapangan. Penulisan di subbab ini harus disusun secara runut dan berkorespondensi langsung untuk menjawab butir-butir Target atau Tujuan Penelitian yang telah Anda tetapkan di Bab I. Gunakan instrumen visual pendukung seperti tabel, diagram batangan, grafik, atau potongan kode/arsitektur visual jika berupa pengembangan sistem. Ingat, fokus di subbab ini adalah \textbf{menyajikan data secara deskriptif dan objektif}, belum masuk ke analisis opini kritis.



% --- Pembahasan ---
\section{Pembahasan}

Bagian ini adalah inti dari orisinalitas pemikiran Anda sebagai peneliti. Di sini Anda wajib melakukan analisis mendalam terhadap hasil penelitian yang ada pada subbab sebelumnya:

\begin{enumerate}
	\item \textbf{Uji Hipotesis (Jika Ada):} Jika penelitian menggunakan metode kuantitatif, deskripsikan proses dan hasil pengujian statistik (seperti uji-t, uji-F, atau analisis jalur) secara transparan untuk membuktikan apakah hipotesis Anda diterima atau ditolak.
	\item \textbf{Sintesis Korelatif:} Hubungkan temuan lapangan Anda dengan kajian teori dan penelitian terdahulu yang relevan (Bab II). Jelaskan apakah hasil penelitian Anda mendukung, memperkuat, atau justru memperlihatkan anomali (kontradiksi) terhadap teori yang sudah mapan.
	\item \textbf{Analisis Kausalitas:} Berikan interpretasi logis mengapa hasilnya bisa demikian dan apa implikasi praktis atau teoretis dari temuan tersebut terhadap pemecahan masalah utama Anda.
\end{enumerate}